% ===== FILE 9/9: Toàn cảnh lịch trình, Chi phí trung thực, Bảng đang chạy, Bảng so sánh, Ưu điểm, Tổng kết, Cảm ơn =====

% --- Slide: Toàn cảnh lịch trình huấn luyện ---
\begin{frame}{Toàn cảnh lịch trình huấn luyện ($30\,000$ iteration)}
\scriptsize
\begin{center}
\begin{tabular}{lll}
\toprule
\textbf{Mốc} & \textbf{Sự kiện} & \textbf{Ghi chú}\\
\midrule
iter $1\to3$ & tăng bậc SH tới $D_{\max}=3$ & mỗi iteration\\
iter $500$ & bắt đầu densify & \texttt{densify\_from\_iter}\\
iter $600,700,\dots$ & clone/split/prune ($\eta$ OR gradient) & mỗi $100$, tới iter $15\,000$\\
iter $3\text{k},6\text{k},9\text{k},12\text{k}$ & reset opacity ($\times0.1$) & đúng $4$ lần, chỉ khi $<15\,000$\\
iter $4\,000$, $8\,000$ & prune cứng thêm ($\alpha<0.1$) & \texttt{prune\_iterations}\\
iter $15\,000$ & \textbf{dừng hẳn} densify/prune/reset & \texttt{densify\_until\_iter}\\
iter $15\,000\to30\,000$ & chỉ còn tối ưu Adam, $N$ không đổi & ---\\
\bottomrule
\end{tabular}
\end{center}
\begin{itemize}\itemsep1pt
    \item \textbf{Không có} cơ chế nào chủ động thay đổi $N$ sau iteration $15\,000$ --- toàn bộ số lượng Gaussian đạt
    được tại mốc đó được giữ nguyên cho tới hết quá trình train, chỉ còn tối ưu giá trị tham số.
    \item Đây là điểm quan trọng khi so sánh 3DGS gốc với SADGS: khác biệt về \textbf{chất lượng/tốc độ hội tụ} chủ yếu
    xảy ra trong $15\,000$ iteration đầu, không phải "số Gaussian cuối cùng khác nhau nhờ tiếp tục tỉa về sau".
\end{itemize}
\end{frame}

% --- Slide: Mô hình chi phí — trung thực ---
\begin{frame}{Chi phí tính toán: mô hình định tính, KHÔNG có số đo thật}
\scriptsize
Repo này \textbf{không chứa} script benchmark thời gian thực (wall-clock) nào đo trực tiếp và so sánh SADGS với 3DGS
gốc trên cùng phần cứng, cùng scene, cùng số lần lặp lại đo. Một mô hình \textbf{định tính} hợp lý (không phải số đo
thực nghiệm) cho chi phí mỗi iteration:
\[
T_{\text{iter}} \approx a\,N \;+\; b\,N\,K \;+\; c\,N\cdot\mathbb{1}[\text{Adam step}] \;+\; F
\]
với $N$ số Gaussian, $K$ số tile trung bình/Gaussian (bị chi phối bởi \texttt{mult} của compact box), $F$ chi phí cố
định mỗi khung hình (forward pass loss, sắp xếp, v.v.).
\begin{itemize}\itemsep1pt
    \item Thành phần \textbf{có thể làm tăng} chi phí/iteration so với đường cơ sở: cập nhật $\eta$ mỗi $10$ iteration
    (thêm một lượt tính structure tensor và chiếu trục), hai optimizer chạy song song trong chế độ \texttt{hybrid}.
    \item Thành phần \textbf{có thể làm giảm} chi phí: \texttt{mult}$=0.7<1$ giảm $K$ so với SnugBox nguyên bản
    (\texttt{mult}$=1$) --- ít tile hơn cho mỗi Gaussian.
    \item \textbf{Nguyên tắc trình bày của bài giảng này:} bất kỳ con số "nhanh hơn $X$ lần", "giây/iteration" nào không
    đi kèm phương pháp đo rõ ràng (phần cứng cụ thể, scene cụ thể, số lần lặp lại đo) sẽ \textbf{không được trích dẫn}
    như một kết quả benchmark đã kiểm chứng.
\end{itemize}
\end{frame}

% --- Slide: Bảng những gì đang thật sự chạy mặc định ---
\begin{frame}{Tóm tắt: những gì ĐANG CHẠY mặc định trong repo này}
\scriptsize
\begin{center}
\begin{tabular}{ll}
\toprule
\textbf{Cơ chế} & \textbf{Vị trí trong code}\\
\midrule
Structure tensor đa tỉ lệ $\to\eta\to$ split/clone/prune & \texttt{utils/freq\_utils.py}, \texttt{train.py:276,341--372}\\
Split dị hướng $k_x k_y k_z$ theo từng trục & \texttt{scene/gaussian\_model.py:638}\\
Tích luỹ nhất quán đa view (high/low ratio) & \texttt{train.py:341--353}\\
Reset opacity $\times0.1$, đúng $4$ lần & \texttt{scene/gaussian\_model.py:415}\\
Trần opacity $\min(\alpha,0.8)$ & \texttt{scene/gaussian\_model.py:1043}\\
Optimizer hybrid (Adam thường $+$ Sparse/Fused Adam) & \texttt{train.py:429--434}\\
Compact box SnugBox, \texttt{mult}$=0.7$ & \texttt{cuda\_rasterizer/auxiliary.h:313}\\
Bộ lọc low-pass 2D, $\kappa=0.1$ & \texttt{cuda\_rasterizer/forward.cu:113}\\
\bottomrule
\end{tabular}
\end{center}
\footnotesize Đây là danh sách đầy đủ những gì bài giảng này đã trình bày chi tiết --- các nhánh tuỳ chọn khác trong
code (nhánh 3DGS gốc dùng làm baseline so sánh ở phần đầu, và một vài cơ chế được thiết kế nhưng hiện đang tắt theo
mặc định) không được đào sâu thêm ở đây vì không ảnh hưởng tới hành vi mặc định của pipeline.
\end{frame}

% --- Slide: Bảng so sánh 3DGS gốc / SADGS (chỉ những khác biệt đang thật sự chạy) ---
\begin{frame}{Bảng so sánh 3DGS gốc vs.\ SADGS (chỉ những khác biệt ĐANG CHẠY)}
\scriptsize
\begin{center}
\begin{tabular}{p{4.6cm}p{3.6cm}p{4.4cm}}
\toprule
\textbf{Tiêu chí} & \textbf{3DGS gốc (baseline)} & \textbf{SADGS (mặc định)}\\
\midrule
Tín hiệu densify chính & gradient vị trí & gradient \textbf{OR} $\eta$ tần số đa tỉ lệ\\
Split dị hướng theo trục & không có (luôn $N{=}2$, mọi trục $/1.6$) & có, $k_{\text{axis}}=\lceil\sqrt{\max(\eta,1)}\rceil$ riêng từng trục\\
Tích luỹ đa view trước quyết định & không (đọc gradient tại 1 lần cuối) & có (\texttt{high\_ratio}/\texttt{low\_ratio}$>0.8$ qua nhiều view)\\
Vị trí Gaussian con khi split & ngẫu nhiên $\mathcal{N}(0,\mathrm{diag}(s)^2)$ & lưới tất định, cách đều $\sqrt{12}\sigma$\\
\bottomrule
\end{tabular}
\end{center}
\footnotesize\textit{Compact box, bộ lọc low-pass 2D và optimizer hybrid là các cơ chế \textbf{chạy giống nhau} ở cả hai
nhánh (xem slide "kế thừa vs mới") --- không đưa vào bảng so sánh vì không phải điểm khác biệt.}
\end{frame}

% --- Slide: Ưu điểm cốt lõi ---
\begin{frame}{Ưu điểm cốt lõi của SADGS: hội tụ đúng chỗ, không phải "ít Gaussian hơn"}
\scriptsize
\begin{itemize}\itemsep2pt
    \item \textbf{Đừng diễn giải sai:} không có bằng chứng trong repo cho thấy SADGS tạo ra \textbf{ít Gaussian hơn} so
    với 3DGS gốc ở cùng mức chất lượng --- cơ chế duy nhất có khả năng chủ động giảm mạnh $N$ ở giai đoạn cuối không
    nằm trong phạm vi những gì đang chạy mặc định (xem slide "đang chạy mặc định").
    \item \textbf{Lợi ích thực sự, có căn cứ trực tiếp từ công thức đã trình bày:}
    \begin{itemize}
        \item Phân bổ ngân sách split \textbf{đúng trục, đúng mức độ} ($k_{\text{axis}}$ tính riêng từng trục, tăng
        theo bậc thang tương xứng với mức độ vi phạm) thay vì luôn chia đôi đều mọi trục như cơ chế gốc.
        \item Quyết định dựa trên \textbf{đa số view đồng thuận} ($>80\%$), giảm hẳn nguy cơ phản ứng theo nhiễu của một
        góc chụp đơn lẻ bất thường.
        \item Tín hiệu $\eta$ đo trực tiếp \textbf{tỉ lệ kích thước Gaussian / chi tiết ảnh}, bổ sung cho gradient (vốn
        chỉ phản ánh trạng thái hội tụ hiện tại, và có thể bị triệt tiêu nội-ảnh giữa các pixel như đã học).
    \end{itemize}
    \item \textbf{Còn bỏ ngỏ, cần nói rõ khi trình bày kết quả:} repo không có số đo PSNR/SSIM/LPIPS hay wall-clock nào
    so sánh trực tiếp hai nhánh trên cùng scene/checkpoint để \textbf{định lượng} các lợi ích lý thuyết vừa nêu.
\end{itemize}
\end{frame}

% --- Slide: Tổng kết ---
\begin{frame}{Tổng kết: từ 3DGS gốc tới SADGS}
\scriptsize
\begin{enumerate}\itemsep2pt
    \item \textbf{Biểu diễn \& render:} Gaussian 3D tường minh, EWA splatting, alpha blending, rasterizer theo tile ---
    nền tảng chung, giữ nguyên giữa hai nhánh.
    \item \textbf{Gradient:} tích luỹ đúng cách (chuẩn trước, cộng sau) tránh triệt tiêu giữa các view; vấn đề triệt
    tiêu thật nằm giữa các pixel trong cùng một ảnh, được xử lý bằng một kênh gradient trị tuyệt đối riêng.
    \item \textbf{ADC gốc:} clone/split/prune dựa thuần trên gradient --- khung sườn chung mà SADGS kế thừa nguyên vẹn.
    \item \textbf{SADGS:} bổ sung tín hiệu $\eta$ (tỉ số kích thước Gaussian / bước sóng texture cục bộ), tích luỹ nhất
    quán qua nhiều view, split dị hướng theo đúng trục vi phạm --- kết hợp \texttt{OR} với gradient chứ không thay thế,
    và đây là phần \textbf{đang chạy mặc định}.
    \item \textbf{Trung thực về hiệu năng:} không có số đo wall-clock/chất lượng nào trong repo để khẳng định "nhanh
    hơn" hay "ít Gaussian hơn" một cách định lượng --- bài giảng này trình bày đúng những gì công thức và code thể
    hiện, không suy diễn thêm ngoài đó.
\end{enumerate}
\end{frame}

% --- Slide: Cảm ơn / Q&A ---
\begin{frame}{Cảm ơn}
\centering
\vspace{1cm}
{\Large Câu hỏi \& thảo luận}\\[0.8cm]
\footnotesize
Toàn bộ công thức trong bài giảng này được đối chiếu trực tiếp với:\\
\texttt{train.py}, \texttt{scene/gaussian\_model.py}, \texttt{utils/freq\_utils.py}, \texttt{utils/loss\_utils.py},\\
\texttt{arguments/\_\_init\_\_.py}, \texttt{cuda\_rasterizer/\{forward.cu, auxiliary.h\}} và \texttt{DOCS/BOOK/}.
\end{frame}
